\section{Discussion}
\label{sec:discussion}

\subsection{What the linter does and does not enforce}

The provenance linter enforces the \emph{syntax} of epistemic hygiene. A paragraph must carry a tag; a \texttt{[SIMULATION\_REHEARSAL]} paragraph must not contain the phrase ``users preferred''; a \texttt{[SOURCE\_CARD:id]} tag must reference a corpus ID that exists. The linter cannot check whether a citation, grounded in a card that exists, actually supports the claim the paragraph makes. The demo guidebook contained exactly this failure: a \texttt{[SOURCE\_CARD:liang\_nejm\_ai\_2024]}-tagged paragraph cited Liang et al.'s 57.4\% helpfulness finding to support a proposition about credibility judgments on AI-authored content \citep{liang2024large}. Liang measured researcher-rated feedback helpfulness on scientific manuscripts. The guidebook used the number to make a claim about a different population evaluating a different kind of content.

A later review wave caught this and the guidebook paragraph was rewritten to accurately describe what Liang measured. The linter did not catch it, and the linter as described in this paper will not catch it on a future run. Citation-misuse detection is an open problem: it requires either a human review step between generation and shipping, or a second agent whose job is to verify that each quoted claim is present in the cited source card's \texttt{allowed\_claims} list. The latter is a natural extension we have not implemented.

\subsection{The \texttt{[AGENT\_INFERENCE]} tag as an escape hatch}

Among the eight provenance tags, seven make a load-bearing claim about where the content came from. The eighth, \texttt{[AGENT\_INFERENCE]}, is a catch-all for reasoning that is neither grounded in a source card nor drawn from the simulated walkthrough. Its force is only as strong as the model's willingness not to route low-confidence content through it.

One concrete failure mode we caught after the first review wave: the model wanted to name outside-corpus adjacent literature (Jakesch, Longoni, Sundar) and chose \texttt{[AGENT\_INFERENCE]} as the wrapper. This is technically honest --- the mention is not grounded in a corpus card --- but it hides the verification requirement from the researcher who reads the guidebook. The fix was to add an \texttt{[UNCITED\_ADJACENT]} tag specifically for this case, with the guidebook-agent prompt updated to treat it as the required wrapper for named outside-corpus papers. The linter accepts either tag; the human-facing distinction is whether the tag commits the researcher to verification.

A stronger version of the system would prohibit \texttt{[AGENT\_INFERENCE]} from containing any named reference that is not in the corpus. We have not enforced this. Doing so would make \texttt{[AGENT\_INFERENCE]} effectively a ``pure reasoning'' tag whose content never references specific external work.

\subsection{Pattern applicability and out-of-domain transfer}

The second-domain benchmark run on the code-review premise exposed a pattern-transfer issue. Two branches were blocked on \texttt{legibility.no\_failure\_signal}, which matches the pattern's trigger heuristic exactly: the proposed tools do not surface when they are wrong. But on one branch, the ``weak override'' pattern fired against a feature whose deliberate purpose was to create friction against engineer over-trust --- the friction was the manipulation under study, not a constraint on user intent. The auditor did not distinguish between ``weak override as unintended constraint'' and ``weak override as intentional manipulation.''

The current mitigation is a note in the \texttt{WORKSHOP\_NOT\_RECOMMENDED} template asking the human reviewer to confirm the pattern fired correctly rather than mis-matched a manipulation-under-test. A stronger version of the system would require each pattern's trigger heuristic to include an explicit applicability check: ``this pattern fires only if the matched constraint is not itself the study's manipulation.'' This is a future-work item \S\ref{sec:limitations}.

\subsection{Why Opus 4.7 for the reviewer stage, given the ablation result}

The ablation found Sonnet 4.6 capable of catching the methodologist confound we had hoped was Opus-specific. This does not mean the routing decision was wrong; it means the justification for the decision changes. Opus produces more precise evidence spans and more specific required-change instructions. On a run where the study lives or dies on a 0.5-second manipulation-check threshold, that precision matters. On a run where the study is earlier in design and the researcher is looking for the overall shape of reviewer objections, Sonnet is substantively equivalent at lower cost.

A defensible policy would be: route initial triage runs to Sonnet, and re-run the reviewer stage on Opus only for branches where the meta-reviewer has returned \texttt{human\_judgment\_required} (meaning the reviewer panel's disagreement is the critical input to the final decision). We have not implemented this tiered policy. The current default is Opus everywhere for the reviewer stage, which the ablation does not support.

\subsection{Rehearsal as pre-registered prior}

The project's tagline is ``rehearsal stage for research; the performance still needs humans.'' The linter's structural job is narrow: simulated walkthroughs may not use evidence language. The more interesting function the tag performs is a productive one. A rehearsal sentence like ``a fast H-key navigator would traverse all three H2s in under four seconds'' is a falsifiable prediction the real study either confirms or refutes. The \texttt{[SIMULATION\_REHEARSAL]} tag reframes simulation-as-evidence into simulation-as-pre-registered-prior.

This reframing matters because a researcher who reads the guidebook and proceeds to a real study will inevitably have their mental model shaped by the rehearsal. The defensive claim that rehearsal is not evidence does not prevent this. What prevents the rehearsal from laundering itself as findings is the requirement that each rehearsal paragraph use hedged language (``would likely,'' ``could plausibly,'' ``a user encountering this might'') and the \texttt{[HUMAN\_REQUIRED]} block at the end of the guidebook. The latter makes explicit the handoff: the researcher must confirm what the rehearsal predicted.
